\documentclass[aspectratio=43,11pt]{beamer}

\usetheme{default}
\usecolortheme{default}
\definecolor{aggiemaroon}{RGB}{80,0,0}
\setbeamercolor{structure}{fg=aggiemaroon}
\setbeamertemplate{navigation symbols}{}\setbeamertemplate{footline}[frame number]
\setbeamertemplate{itemize items}[circle]

\usepackage{amsmath,amssymb,booktabs}

%\title{A Simple Manual Example of MCMC}
%\subtitle{Metropolis-Hastings for a Bernoulli Parameter}
%\author{}
%\date{}

\title{Decision and Risk Analysis}
\subtitle{Lecture 7: \textbf{A Simple Manual Example of MCMC}\\
Metropolis-Hastings for a Bernoulli Parameter}
\author{}
\date{
\color{aggiemaroon}{
DAEN 427 / ISEN 427 / ISEN 627 - Fall 2026\\
\textbf{Alexander Abuabara}
}
}
\titlegraphic{\vfill \includegraphics[height=1.2cm]{TAMU_logo.png}}

\begin{document}

\begin{frame}[plain]
  \titlepage
\end{frame}

\begin{frame}{}
We want to understand how MCMC can generate samples from a posterior distribution.

\vspace{0.4cm}
Suppose we observe four Bernoulli trials:
\[
 y = (1,1,1,0)
\]

That is:
\[
 3 \text{ successes}, \qquad 1 \text{ failure}
\]

\vspace{0.3cm}
The unknown parameter is:
\[
\theta = \Pr(y=1)
\]
\end{frame}

\begin{frame}{Prior and Likelihood}
Assume a uniform prior:
\[
\theta \sim \text{Uniform}(0,1)
\]

The likelihood for 3 successes and 1 failure is:
\[
 p(y \mid \theta) \propto \theta^3(1-\theta)^1
\]

Since the prior is uniform, the posterior is proportional to:
\[
 p(\theta \mid y) \propto \theta^3(1-\theta)
\]

We define the unnormalized posterior density:
\[
 f(\theta)=\theta^3(1-\theta)
\]
\end{frame}

\begin{frame}{Metropolis-Hastings Rule}
Start at some current value:
\[
\theta_t
\]

Propose a new value:
\[
\theta^*
\]

Compute the acceptance ratio:
\[
 r = \frac{f(\theta^*)}{f(\theta_t)}
\]

Then compute:
\[
 \alpha = \min(1,r)
\]

Accept the proposed value if:
\[
 u < \alpha, \qquad u \sim \text{Uniform}(0,1)
\]
\end{frame}

\begin{frame}{Step 1: Move from 0.50 to 0.60}
Start at:
\[
\theta_0 = 0.50
\]

Compute:
\[
 f(0.50)=0.50^3(1-0.50)=0.125 \times 0.50=0.0625
\]

Propose:
\[
\theta^*=0.60
\]

Compute:
\[
 f(0.60)=0.60^3(1-0.60)=0.216 \times 0.40=0.0864
\]

Acceptance ratio:
\[
 r=\frac{0.0864}{0.0625}=1.3824
\]

So:
\[
\alpha=1
\]

Accept:
\[
0.50 \rightarrow 0.60
\]
\end{frame}

\begin{frame}{Step 2: Move from 0.60 to 0.80}
Current value:
\[
\theta_t=0.60
\]

Propose:
\[
\theta^*=0.80
\]

Compute:
\[
 f(0.80)=0.80^3(1-0.80)=0.512 \times 0.20=0.1024
\]

Acceptance ratio:
\[
 r=\frac{0.1024}{0.0864}=1.185
\]

Thus:
\[
\alpha=1
\]

Accept:
\[
0.60 \rightarrow 0.80
\]
\end{frame}

\begin{frame}{Step 3: A Rejected Proposal}
\small
Current value:
\[
\theta_t=0.80
\]

Propose:
\[
\theta^*=0.95
\]

Compute:
\[
 f(0.95)=0.95^3(1-0.95)\approx 0.8574\times 0.05=0.0429
\]

Acceptance ratio:
\[
 r=\frac{0.0429}{0.1024}\approx 0.419
\]

So:
\[
\alpha=0.419
\]

Suppose we draw: % {\tiny(random value between $0-1$)}:
\[
u=0.70
\]

Since:
\[
0.70>0.419
\]

reject the proposal and stay at:
\[
\theta=0.80
\]
\end{frame}

\begin{frame}{Step 4: Move from 0.80 to 0.70}
Current value:
\[
\theta_t=0.80
\]

Propose:
\[
\theta^*=0.70
\]

Compute:
\[
 f(0.70)=0.70^3(1-0.70)=0.343 \times 0.30=0.1029
\]

Acceptance ratio:
\[
 r=\frac{0.1029}{0.1024}\approx1.005
\]

Thus:
\[
\alpha=1
\]

Accept:
\[
0.80 \rightarrow 0.70
\]
\end{frame}

\begin{frame}{Manual Chain}
The chain we recorded is:
\[
0.50,\;0.60,\;0.80,\;0.80,\;0.70
\]

\vspace{0.4cm}
Notice the repeated value:
\[
0.80,\;0.80
\]

This happens because a rejected proposal keeps the chain at the current state.

\vspace{0.4cm}
After many iterations, the histogram of the chain approximates the posterior:
\[
p(\theta \mid y)
\]
\end{frame}

\begin{frame}{Summary Table}
\centering
\small
\begin{tabular}{ccccccc}
\toprule
Step & Current & Proposal & $f$(Current) & $f$(Proposal) & $\alpha$ & Result \\
\midrule
1 & 0.50 & 0.60 & 0.0625 & 0.0864 & 1.000 & Accept \\
2 & 0.60 & 0.80 & 0.0864 & 0.1024 & 1.000 & Accept \\
3 & 0.80 & 0.95 & 0.1024 & 0.0429 & 0.419 & Reject \\
4 & 0.80 & 0.70 & 0.1024 & 0.1029 & 1.000 & Accept \\
\bottomrule
\end{tabular}
\end{frame}

\begin{frame}{Key Idea}
MCMC does not only search for the best value.

\vspace{0.4cm}
It constructs a chain whose long-run behavior matches the posterior distribution.

\vspace{0.4cm}
Regions with higher posterior probability are visited more often.

\vspace{0.4cm}
Regions with lower posterior probability are visited less often, but not necessarily never.

\vspace{0.4cm}
That is why MCMC provides uncertainty, not just a point estimate.
\end{frame}

\begin{frame}{} % Connection to Bayesian Regression
In the simple example, the unknown was one parameter:
\[
\theta
\]

In Bayesian linear regression, the unknowns are usually:
\[
(\beta_0,\beta_1,\sigma)
\]

MCMC samples from:
\[
p(\beta_0,\beta_1,\sigma \mid x,y)
\]

Each draw gives one plausible regression model:
\[
y = \beta_0 + \beta_1 x + \varepsilon
\]

The collection of draws gives the full posterior uncertainty.
\end{frame}

\end{document}